\section{Hardness results}
\label{sec:hardness}

First, we consider the NP-hardness of the following problem.

\begin{problem}{\probname{$1$-Tight OCC Detection}}
  \Input & A graph $G$.\\
  \Prob & Output a non-empty $1$-tight OCC $(X_B,X_C,X_R)$ in $G$ or conclude that no non-empty $1$-tight OCC exists.
\end{problem}

To prove that this problem is NP-hard under Turing reductions,
we consider the following version of \PrbOCT.

\begin{problem}{\probname{Lower-Bounded OCT}}
  \Input & A graph $G$ and $k$ vertex-disjoint odd cycles $\{C_i\}$ in $G$.\\
  \Prob & Is there a set $S \subseteq V(G)$ of size at most $k$ such that $G-S$ is bipartite?
\end{problem}

We claim that \probname{$1$-Tight OCC Detection} is as hard as \probname{Lower-Bounded OCT}.

% \bmp{The exposition in this section overlooks the distinction between NP-hardness under Turing reductions versus under Karp reductions. Typically, when one sees '$X$ is NP-hard', it means NP-hard under Karp reductions: for each problem~$Y$ in NP, there is a polynomial-time reduction that maps each instance of~$Y$ to a single instance of~$X$ with the same YES/NO answer. That is not the type of NP-hardness we get here, since we have to make \emph{multiple} calls to an algorithm for 1-tight OCC detection to solve lower-bounded OCT. Our proof effectively shows hardness under Turing reductions, i.e., if there was a poly-time algorithm for 1-tight OCC detection, then there would be a polynomial-time algorithm for all problems in NP. I suggest staying closer to the formulation from the journal version of the previous paper, or explicitly mentioning 'NP-hard under Turing reductions', to make this clear.}

\begin{lemma}\label{lem:1-tight-occ-reduction}
  Assuming $\textnormal{\textsf{P}} \neq \textnormal{\textsf{NP}}$,
  if there is no polynomial-time algorithm for \probname{Lower-Bounded OCT}, then
  there is no polynomial-time algorithm for \probname{$1$-Tight OCC Detection}.
\end{lemma}

\begin{proof}
  Assume \probname{$1$-Tight OCC Detection} is solvable in polynomial time.
  Then, consider the following algorithm for \probname{Lower-Bounded OCT}.

  \begin{enumerate}
    \item Let $C \gets \{C_i\}$. Initially, $|C|=k$.
    \item Return \textsf{Yes} if $C=\emptyset$.
    \item Run an algorithm for \probname{$1$-Tight OCC Detection} on $G$.
    If it finds a non-empty $1$-tight OCC $(X_B,X_C,X_R)$, then
    let $G \gets X_R$ and $C \gets \{C_i \in C \mid X_C \cap V(C_i) \neq \emptyset\}$,
    i.e. remove odd cycles in $C$ that hit any vertex in $X_C$,
    and then continue to Step 2.
    %
    Otherwise, return \textsf{No}.
  \end{enumerate}

  Let $S$ be the union of $X_C$ found in Step 3.
  If the algorithm returns \textsf{Yes}, then
  it is clear to see that $|S|=k$ and $G-S$ is bipartite.
  $(G,\{C_i\})$ is a yes-instance for \probname{Lower-Bounded OCT}.
%
  If the algorithm returns \textsf{No}, then
  there exists an induced subgraph $G'$ of $G$ such that
  for any OCT $S$ of $G'$, there exists a cycle $C_i$ such that $|V(C_i) \cap S|>1$.
  Hence, $\oct(G)>k$, and $(G,\{C_i\})$ is a no-instance.
%
  If \probname{$1$-Tight OCC Detection} is poly-time solvable,
  then so is \probname{Lower-Bounded OCT}.
\end{proof}

Now, let us show that \probname{Lower-Bounded OCT} is actually NP-hard.
We adopt an argument from \cite[Lemma 16]{DonkersJ19} but with simpler gadgets.

\begin{lemma}\label{lem:lb-oct-reduction}
  \probname{Lower-Bounded OCT} is NP-hard.
\end{lemma}

\begin{proof}
  We reduce from \probname{3-SAT}, a version of \probname{Satisfiability (SAT)}
  where every clause includes exactly $3$ literals, to \probname{Lower-Bounded OCT}.
  %
  Let $X=x_1,\ldots,x_n$ be the variables and $\Phi=\phi_1,\ldots,\phi_m$ be the clauses
  appearing in a SAT instance.

  We construct a graph $G$ as follows:
  \begin{itemize}
    \item For each variable $x_i$, create a new triangle and let $v_i$ and $\overline{v_i}$
    be two of the triangle's vertices.
    These vertices represent SAT literals $x_i$ and $\overline{x_i}$ as shown in \Cref{fig:proof-np-hardness} (left).
    \item For each clause $\phi_j=\ell_{j1} \lor \ell_{j2} \lor \ell_{j3}$, introduce a gadget illustrated in \Cref{fig:proof-np-hardness} (right).
    We then identify vertices $s_1, s_2, s_3$ and the vertices representing $\ell_{j1}$, $\ell_{j2}$, $\ell_{j3}$,
    respectively.
  \end{itemize}

  It is clear to see that this construction can be done in polynomial time.
  Also observe that there are $n$ disjoint triangles for SAT variables
  and $2m$ disjoint triangles for clause gadgets.
  %
  By letting $k=n+2m$, we have that $G$ contains $k$ vertex-disjoint triangles, which are odd cycles.
  We claim that $(X,\Phi)$ is satisfiable if and only if $G$ contains an OCT of size at most $k$.

  First, assume $(X,\Phi)$ is satisfiable.
  Let $U \subseteq V(G)$ be the vertices that correspond to a certificate,
  and consider all odd cycles present in $G-U$.
  %
  For each variable gadget, we know that either $x_i$ or $\overline{x_i}$ is in $U$,
  so no triangles remain in $G-U$.
  %
  For each clause gadget, we know that at least one of $s_1, s_2, s_3$ is in $U$.
  In any cases, we can choose $2$ vertices hitting all remaining triangles.
  %
  Hence, $\oct(G-U) \leq 2m$ and $\oct(G) \leq n + 2m = k$.

  Conversely, assume $G$ has an OCT $S \subseteq V(G)$ of size at most $k$.
  Since $G$ has $k$ vertex-disjoint triangles, $S$ contains exactly one vertex from each of them.
  For every variable gadget, exactly one vertex is in $S$.
  For every clause gadget excluding $s_1$, $s_2$, $s_3$, exactly $2$ vertices must be in $S$,
  which leads to the fact that at least one of $s_1$, $s_2$, $s_3$ is in $S$.
  %
  Finally, we construct a solution for \probname{3-SAT}.
  For each $1 \leq i \leq n$,
  we set $x_i$ to \textsf{true} if $v_i \in S$,
  and \textsf{false} otherwise.
  This by construction forms a certificate, and $(X,\Phi)$ is satisfiable.
\end{proof}

Combining \Cref{lem:1-tight-occ-reduction} and \Cref{lem:lb-oct-reduction},
we obtain the following result.

\begin{theorem} \label{thm:1tight:occ:nphard}
  \probname{$1$-Tight OCC Detection} is NP-hard under Turing reductions.
\end{theorem}

Next, we show that finding a non-empty 1-tight OCC of bounded width is W[1]-hard,
parameterized by the width.

\begin{problem}{\probname{Bounded-Width $1$-Tight OCC Detection}}
  \Input & A graph $G$ and an integer $k$.\\
  Parameter: & $k$. \\
  \Prob & Does $G$ admit a non-empty $1$-tight OCC $(X_B,X_C,X_R)$ with $|X_C|\leq k$?
\end{problem}

We will show the W[1]-hardness of this problem by a parameterized reduction from
the well-known \probname{Multicolored Clique} problem.

\begin{problem}{\probname{Multicolored Clique}}
  \Input & A graph $G$, an integer $k$, and a partition of $V(G)$ into sets $V_1,\ldots,V_k$.\\
  Parameter: & $k$. \\
  \Prob & Does $G$ include as a subgraph such a clique $S$ such that
  for each $1 \leq i \leq k$ we have $|S \cap V_i|=1$?
\end{problem}

\begin{figure*}[t]
  \centering
  \tikzstyle{large} = [circle, fill=white, text=black, draw, thick, scale=1, minimum size=0.8cm, inner sep=1.5pt]
  \tikzstyle{plain} = [circle, fill=white, text=black, draw, thick, scale=1, minimum size=0.4cm, inner sep=1.5pt]
  \tikzstyle{small} = [circle, fill=white, text=black, draw, thick, scale=1, minimum size=0.2cm, inner sep=1.5pt]
  \tikzstyle{black} = [circle, fill=gray, text=black, draw, thick, scale=1, minimum size=0.2cm, inner sep=1.5pt]
  \tikzmath{\xunit = 0.6; \yunit = 1.0; }

  \begin{minipage}[m]{.30\textwidth}
    \vspace{0pt}
    \centering
    \begin{tikzpicture}
        % vertices
        \node[plain,label=below:$v_i$] (a2) at (2 * \xunit, 1 * \yunit) {};
        \node[plain,label=below:$\overline{v_i}$] (a3) at (4 * \xunit, 1 * \yunit) {};
        \node[plain] (b1) at (3 * \xunit, 2 * \yunit) {};
        
        \draw (b1)--(a2)--(a3)--(b1);
    \end{tikzpicture}
  \end{minipage}
  \hfill
  %
  \begin{minipage}[m]{.66\textwidth}
      \vspace{0pt}
      \centering
      \begin{tikzpicture}
          % vertices
          \node[plain,dotted,label=below:$s_{1}$] (s1) at (1 * \xunit, 0) {};
          \node[plain,dotted,label=below:$s_{2}$] (s2) at (5 * \xunit, 0) {};
          \node[plain,dotted,label=below:$s_{3}$] (s3) at (9 * \xunit, 0) {};
          \node[plain] (a1) at (0 * \xunit, 1 * \yunit) {};
          \node[plain] (a2) at (2 * \xunit, 1 * \yunit) {};
          \node[plain] (a3) at (4 * \xunit, 1 * \yunit) {};
          \node[plain] (a4) at (6 * \xunit, 1 * \yunit) {};
          \node[plain] (a5) at (8 * \xunit, 1 * \yunit) {};
          \node[plain] (a6) at (10 * \xunit, 1 * \yunit) {};
          \node[plain] (b1) at (3 * \xunit, 2 * \yunit) {};
          \node[plain] (b2) at (7 * \xunit, 2 * \yunit) {};
          
          \draw (a1)--(a2)--(a3)--(a4)--(a5)--(a6);
          \draw (a1)--(s1)--(a2);
          \draw (a3)--(s2)--(a4);
          \draw (a5)--(s3)--(a6);
          \draw (a2)--(b1)--(a3);
          \draw (a4)--(b2)--(a5);
          
          % \draw (u23)--(u32);
          % \draw (u31)--(u32);
      \end{tikzpicture}
  \end{minipage}
%
  \caption{%
  Gadgets used for the proof of \Cref{lem:lb-oct-reduction}.
  The left shows a \emph{variable gadget}, including two vertices representing literals.
  The right shows a \emph{clause gadget}, where $s_1$, $s_2$, $s_3$ are connected to
  corresponding literals.
  }
  \label{fig:proof-np-hardness}
\end{figure*}


\begin{theorem} \label{thm:tightocc:w1hard}
  \probname{Bounded-Width $1$-Tight OCC Detection} is W[1]-hard.
\end{theorem}

\begin{proof}
  Consider an instance $(G,\chi,k)$ of \probname{Multicolored Clique},
  where $\chi\colon V(G) \to [k]$ denotes the \textit{color} of each vertex such that
  $\chi(v)=\ell$ if and only if $v \in V_\ell$.
%
  Without loss of generality, we assume $k \geq 2$ and $n > k + 2$.
%
  Then, we construct an input $(G',k')$ for \probname{Bounded-Width $1$-Tight OCC Detection} as follows
  (see \Cref{fig:proof-w1-hardness} for an illustration).
%
  \begin{enumerate}
    \item For each vertex $i \in V(G)$, introduce ($k-1$) vertices $U_{i}:=\{u_{i\ell} \mid \ell \in [k] \setminus \{\chi(i)\}\}$ to $G'$.
    \item For each edge $ij \in E(G)$, add $4$ vertices $x_{ij}$, $x_{ij}'$, $y_{ij}$, and $y_{ij}'$ to $G'$.
    \item Insert edges to turn $\bigcup_{i \in V(G)}U_i$ into a clique in $G'$.
    \item For each edge $ij \in E(G)$, add edges from $U_i \cup U_j$ to $x_{ij}$,
    $x_{ij}$ to $x_{ij}'$, and $x_{ij}'$ to $u_{i \chi(j)}$.
    Similarly, add edges from $U_i \cup U_j$ to $y_{ij}$,
    $y_{ij}$ to $y_{ij}'$, and $y_{ij}'$ to $u_{j \chi(i)}$.
  \end{enumerate}

  Let $k'=k(k-1)$, and we claim that $(G',k')$ for \probname{Bounded-Width $1$-Tight OCC Detection}
  \probname{Multicolored Clique}  is an equivalent instance to $(G,\chi,k)$ for \probname{Multicolored Clique}.

  \begin{claim}\label{clm:w1-hardness-forward}
    If $G$ has a multicolored clique of size $k$, then $G'$ has a non-empty $1$-tight OCC of width $k'$.
  \end{claim}

  \begin{claimproof}
    Suppose $S$ is a multicolored clique of size $k$ in $G$.
    Define a $1$-tight OCC $(X_B,X_C,X_R)$ as follows:
    \begin{itemize}
      \item $X_B = \bigcup_{i,j \in S:i \neq j} \{ x_{ij}, x_{ij}', y_{ij}, y_{ij}' \}$
      \item $X_C = \bigcup_{i \in S} U_i$
      \item $X_R = V(G') \setminus (X_B \cup X_C)$
    \end{itemize}

    It is clear to see that $|X_C|=k(k-1)=k'$.
    Also, by construction, $N(X_C)=\bigcup_{i,j \in S \mid i \neq j} (U_i \cup U_j \cup \{ u_{i \chi (j)}, u_{j \chi (i)} \}) = \bigcup_{i \in S} U_i = X_C$,
    which remains us to show that $G'[X_B \cup X_C]$ contains $k'$ vertex-disjoint odd cycles.
%
    Notice that $G'[X_B \cup X_C]$ contains the set of triangles
    $T = \bigcup_{i,j \in S \mid i \neq j} \{u_{i \chi(j)}x_{ij}x_{ij}', u_{j \chi(i)}y_{ij}y_{ij}'\}$.
    Since $|T|=\binom{k}{2}\cdot 2=k'$ and $x_{ij},x_{ij}',y_{ij},y_{ij}'$ are unique in $T$,
    it suffices to show that $\left|\bigcup_{i,j\in S \mid i\neq j} \{u_{i\chi(j)}, u_{j\chi(i)}\}\right|=k'$.
%
    Since $S$ is a multicolored clique, for every pairwise distinct $i,j\in S$, we have $\chi(i) \neq \chi(j)$.
    Hence, $\bigcup_{i,j\in S \mid i\neq j} \{u_{i\chi(j)}, u_{j\chi(i)}\}\supseteq
    \bigcup_{i\in S} \left(\bigcup_{j \in S \setminus \{i\}} u_{i \chi(j)}\right)
    =\bigcup_{i\in S} \bigcup_{\ell \in [k] \setminus \{\chi(i)\}}u_{i\ell}
    =\bigcup_{i\in S} U_i$.
    $T$ indeed consists of $\left| \bigcup_{i\in S} U_i \right|=k(k-1)=k'$ vertex-disjoint triangles,
    and $(X_B,X_C,X_R)$ is a $1$-tight OCC.
  \end{claimproof}

  Finally, we show the backward implication.

  \begin{claim}\label{clm:w1-hardness-backward}
    If $G'$ has a non-empty $1$-tight OCC of size $k'$, then $G$ has a multicolored clique of size $k$.
  \end{claim}

  \begin{claimproof}
    Let $(X_B,X_C,X_R)$ be a non-empty $1$-tight OCC of size $k'$ in $G'$.
    We write $U$ for $\bigcup_{i\in S}U_i$ and $W$ for $\bigcup_{e \in E(G)} \{x_e, x_e', y_e, y_e'\}$.

    First, we show that $X_B \subseteq W$.
    Assume not.
    Then, there exists a vertex $u \in X_B \cap U$.
    Since $X_B$ is bipartite and $U$ forms a clique in $G'$,
    we have $|X_B \cap U| \leq 2$.
    %
    We know that $N_{G'}(u) \subseteq X_B \cup X_C$,
    and thus $|N_{G'}(u) \setminus X_B| \leq |X_C|$.
    We have:
    $|N_{G'}(u) \setminus X_B|\geq |N_{G'}(u)|-1
    \geq (|U|-1)-1 = n(k-1)-2 > ((k+2)(k-1)-2) = k^2 + k - 4 \geq k^2-k=k'$.
    %
    This implies $|X_C|>k'$, a contradiction.

    Next, we show that $X_C \subseteq U$.
    For the sake of contradiction, assume there exists a vertex $x \in X_C \cap W$.
    Then, there must be an odd cycle in $G'[\{x\} \cup X_B]$, but because
    $\{x\} \cup X_B \subseteq W$ is bipartite, we reach a contradiction.

    Now that we know $X_B \subseteq W$ and $X_C \subseteq U$,
    for each $u \in X_B$, there is an odd cycle in $G'[\{u\} \cup X_B]$.
    This happens only when $u x_e x_e'$ or $u y_e y_e'$ forms a triangle for some $e \in E(G)$,
    which leads to $|X_B| \geq 2|X_C|$.

    Let $E' = \{e \in E(G) \mid \{x_e, y_e\} \cap X_B \neq \emptyset\}$.
    If $ij \in E'$, then all the vertices in $C_i \cup C_j$ are in $X_C$,
    so we can write $X_C=\bigcup_{i \in V'}C_i$ for some $V' \subseteq V(G)$.
    Furthermore, from $|X_C|=k'=k(k-1)$, we know $|V'|=k$.
    %
    From $E' \subseteq \binom{V'}{2}$, we have $|X_B| \leq 4|E'| \leq 2|V'|(|V'|-1)=2k'$.
    On the other hand, $|X_B| \geq 2|X_C| = 2k'$.
    These inequalities hold only when $E'=\binom{V'}{2}$.

    We claim that $V'$ is a multicolored clique in $G$.
    Let $ij$ be an edge in $E'$.
    By construction, triangles $u_{i \chi(j)}x_{ij}x_{ij}'$ and $u_{j \chi(i)}y_{ij}y_{ij}'$
    are in $G'[X_B \cup X_C]$, and thus $\chi(i)\neq \chi(j)$.
    Since $V'$ forms a clique in $G$ such that any vertex pair has different colors,
    $V'$ is a multicolored clique.
  \end{claimproof}

  From Claims \ref{clm:w1-hardness-forward} and \ref{clm:w1-hardness-backward},
  we conclude that $(G,\chi,k)$ and $(G',k')$ are equivalent,
  and \probname{Bounded-Width $1$-Tight OCC Detection} is W[1]-hard.
\end{proof}

\begin{figure*}[t]
  \centering
  \tikzstyle{large} = [circle, fill=white, text=black, draw, thick, scale=1, minimum size=0.8cm, inner sep=1.5pt]
  \tikzstyle{plain} = [circle, fill=white, text=black, draw, thick, scale=1, minimum size=0.5cm, inner sep=1.5pt]
  \tikzstyle{small} = [circle, fill=white, text=black, draw, thick, scale=1, minimum size=0.2cm, inner sep=1.5pt]
  \tikzstyle{black} = [circle, fill=gray, text=black, draw, thick, scale=1, minimum size=0.2cm, inner sep=1.5pt]
  \tikzmath{\xunit = 0.8; \yunit = 0.8; \dx = 1.0; \dy = 0.4; \cx=0.8; }

  \begin{minipage}[m]{.96\textwidth}
      \vspace{0pt}
      \centering
      \begin{tikzpicture}
          % vertices
          \node[small,label=below:$u_{12}$] (u12) at (0 + 1 * \dx, 1 - 1 * \dy) {};
          \node[small,label=below:$u_{13}$] (u13) at (0 + 2 * \dx, 1 - 2 * \dy) {};

          \node[small,label={[shift={(-0.5, -0.1)}]$u_{21}$}] (u21) at (4, 1) {};
          \node[small,label=below:$u_{23}$] (u23) at (4 + 2 * \dx, 1 - 2 * \dy) {};

          \node[small,label={[shift={(-0.5, -0.1)}]$u_{31}$}] (u31) at (8, 1) {};
          \node[small,label=below:$u_{32}$] (u32) at (8 + 1 * \dx, 1 - 1 * \dy) {};

          \node[small,label=above:$x_{12}$] (x12) at (2.6, 3) {};
          \node[small,label=above:$x_{12}'$] (x12p) at (2.6 - \cx, 3.2) {};
          \node[small,label=above:$y_{12}$] (y12) at (3.4, 3) {};
          \node[small,label=above:$y_{12}'$] (y12p) at (3.4 + \cx, 3.2) {};

          \node[small,label=above:$x_{23}$] (x23) at (6.6, 3) {};
          \node[small,label=above:$x_{23}'$] (x23p) at (6.6 - \cx, 3.2) {};
          \node[small,label=above:$y_{23}$] (y23) at (7.4, 3) {};
          \node[small,label=above:$y_{23}'$] (y23p) at (7.4 + \cx, 3.2) {};

          %edges
          \draw (u12)--(x12);
          \draw (u13)--(x12);
          \draw (u21)--(x12);
          \draw (u23)--(x12);
          \draw (u12)--(y12);
          \draw (u13)--(y12);
          \draw (u21)--(y12);
          \draw (u23)--(y12);

          \draw (x12)--(x12p);
          \draw (y12)--(y12p);
          \draw (x23)--(x23p);
          \draw (y23)--(y23p);

          \draw (x12p)--(u12);
          \draw (y12p)--(u21);

          \draw (u21)--(x23);
          \draw (u23)--(x23);
          \draw (u31)--(x23);
          \draw (u32)--(x23);
          \draw (u21)--(y23);
          \draw (u23)--(y23);
          \draw (u31)--(y23);
          \draw (u32)--(y23);

          \draw (x23p)--(u23);
          \draw (y23p)--(u32);

          \draw (u12)--(u13);
          \draw (u12)--(u21);
          \draw (u12)--(u23);
          \draw (u12)--(u31);
          \draw (u12)--(u32);
          \draw (u13)--(u21);
          \draw (u13)--(u23);
          \draw (u13)--(u31);
          \draw (u13)--(u32);
          \draw (u21)--(u23);
          \draw (u21)--(u31);
          \draw (u21)--(u32);
          \draw (u23)--(u31);
          \draw (u23)--(u32);
          \draw (u31)--(u32);

          \node[] at (9, 3.1) {\large $\cdots$};
          \node[] at (11, 0.4) {\large $\cdots$};


          \draw[dashed,rounded corners] (0, 1.5) rectangle (-0.5 + 3 * \dx, 1 - 3.6 * \dy);
          \draw[dashed,rounded corners] (3.1, 1.5) rectangle (3.5 + 3 * \dx, 1 - 3.6 * \dy);
          \draw[dashed,rounded corners] (7.1, 1.5) rectangle (7.5 + 3 * \dx, 1 - 3.6 * \dy);
          \node[] at (0.4, -0.2) {$U_1$};
          \node[] at (3.5, -0.2) {$U_2$};
          \node[] at (7.5, -0.2) {$U_3$};

        \draw [decorate,decoration={brace,amplitude=5pt}]
        (11.5, 1.5) -- (11.5, -0.5) node[midway,xshift=20pt]{$U$};

        \draw [decorate,decoration={brace,amplitude=5pt}]
        (11.5, 3.5) -- (11.5, 2.5) node[midway,xshift=20pt]{$W$};
      \end{tikzpicture}
  \end{minipage}
%
  \caption{%
  A visualization of an auxiliary graph constructed from an instance $(G,k)$ of \probname{Multicolored Clique}.
  The set $U$ consists of $n(k-1)$ vertices representing the vertices in $G$ and
  their adjacent colors.
  The set $W$ consists of $4m$ vertices representing the edges in $G$ replaced by a gadget for 
  creating odd cycles.
  }
  \label{fig:proof-w1-hardness}
\end{figure*}

